
En
effet, traitons d'abord le cas @@M0@@. Soient @@M1@@ pour @@M2@@ et @@M3@@. Alors (iii) est évident et (ii) est 2.7 (i).

Soient @@M4@@ dans @@M5@@ et @@M6@@. Alors @@M7@@ et @@M8@@ est inversible en dehors de l'ensemble fini @@M9@@ des points singuliers de @@M10@@ (voir FGA 149 \S 6). Soit @@M11@@ une section hyperplane normale de @@M12@@, qui ne contient pas @@M13@@. Alors @@M14@@ et on a une suite exacte :
@@M15@@
Or, @@M16@@ pour tout @@M17@@, et @@M18@@ pour @@M19@@ par 2.5. Donc @@M20@@ est 3-régulier et par 1.4, @@M21@@ est @@M22@@-régulier. Par conséquent, @@M23@@ pour @@M24@@\,; d'où (i).

Supposons maintenant que @@M25@@ et raisonnons par récurrence sur @@M26@@. Soit @@M27@@ une section hyperplane normale de @@M28@@ et considérons la suite exacte :
@@M29@@
Par 1.7, @@M30@@. Donc pour @@M31@@, l'hypothèse de récurrence donne @@M32@@, et comme @@M33@@ par (ii) pour @@M34@@, 2.4 donne @@M35@@\,; d'où @@M36@@. Mais, pour @@M37@@, @@M38@@ (voir FGA 149 \S 6)\,; d'où (i).

Comme hypothèse de récurrence, on suppose que les polynômes @@M39@@ et @@M40@@ sont définis pour @@M41@@. Alors pour tout @@M42@@, @@M43@@ grâce à 2.9.1. Donc par (i) 
